\documentclass[11pt]{article}
\usepackage[a4paper,margin=2cm]{geometry}
\usepackage{xcolor}
\usepackage{enumitem}
\usepackage{parskip}
\usepackage{sectsty}
\usepackage{titlesec}

\definecolor{teal}{HTML}{0E7C66}
\definecolor{amber}{HTML}{B5651D}
\definecolor{slate}{HTML}{334155}
\definecolor{rule}{HTML}{99C9BE}

\titleformat{\section}{\large\bfseries\color{teal}}{}{0pt}{}
\titlespacing{\section}{0pt}{14pt}{4pt}

\newcommand{\slide}[3]{%
  \section*{\color{teal}\rule[2pt]{14pt}{2pt}\ \ Slide #1 \,\textbullet\, #2 \hfill {\normalsize\color{slate}#3}}%
}
\newcommand{\cue}[1]{{\color{amber}\textbf{On screen:} #1}\par\vspace{2pt}}

\begin{document}

\begin{center}
  {\LARGE\bfseries\color{teal} Presentation Script}\\[4pt]
  {\large How Far Can a PSRAM-less ESP32 Go?}\\[2pt]
  {\normalsize A TinyML Deployment Study for WiFi CSI Human Activity Recognition}\\[6pt]
  {\small Muhammad Ahmad \,\textbullet\, L1F22BSCS0634 \,\textbullet\, Section H9 \,\textbullet\, University of Central Punjab}\\[2pt]
  {\small Target length: about 9--10 minutes \,\textbullet\, 12 slides}
\end{center}

\hrule
\vspace{6pt}

{\small\itshape Tips: speak slowly, pause after each key number, and look up at the audience on the bold lines. Advance slides with the right-arrow or space bar. The numbers in this script match the slides exactly.}

\slide{1}{Title}{$\approx$25s}
\cue{Title, your name, the radiating 3D WiFi scene.}
Good morning, everyone. My name is Muhammad Ahmad, and today I will present my research: \textit{``How Far Can a PSRAM-less ESP32 Go?''} It is a TinyML deployment study for WiFi-based human activity recognition. The question behind the whole project is simple: can the cheapest WiFi chip you can buy, a roughly three-dollar ESP32, run modern activity-recognition AI on its own, with no help from a server?

\slide{2}{Why WiFi sensing}{$\approx$40s}
\cue{Three cards: no camera, no wearable, already deployed.}
First, why use WiFi for this at all? When a person moves, their body bends and scatters the WiFi signal that is already in the room. A receiver can read those tiny changes and tell what the person is doing, with no camera and nothing worn on the body. That makes it ideal for privacy-sensitive places like homes, hospitals, and elderly care, where a camera would be intrusive or simply not allowed. And because it reuses WiFi that is already everywhere, it can scale very cheaply.

\slide{3}{What is WiFi CSI HAR}{$\approx$45s}
\cue{Live CSI spectrogram on the left, four numbered steps on the right.}
So what exactly are we sensing? The signal we use is called Channel State Information, or CSI. The idea is this: the router constantly transmits packets; a moving body perturbs them; CSI records how the signal amplitude changes across fifty-two subcarriers over time; and a machine-learning model maps that pattern to an activity such as walking, sitting, or falling. On the left you can see a live CSI spectrogram. That bright band moving through it is, in effect, a body moving through the channel.

\slide{4}{Live demo (interactive)}{$\approx$50s}
\cue{Interactive 3D scene. Drag to orbit; click the scenario buttons.}
To make this concrete, here is an interactive view. \textbf{[Drag to rotate the scene.]} This is the WiFi field radiating from the sensor, and the green figure is a detected person. Let me change scenarios. \textbf{[Click ``Walking''.]} Here the person is walking, and the floor heatmap follows them. \textbf{[Click ``Fall''.]} And here a fall is detected, and the system flags it immediately. I want to be honest that this is an illustration of the concept, not my exact classifier output, but it shows what WiFi presence sensing actually looks like.

\slide{5}{The gap}{$\approx$40s}
\cue{Prior work vs this work, and the real ESP32 board photo.}
Now, here is the gap I found in the literature. Almost everyone who runs these models on a chip uses the more powerful ESP32-S3, which has extra PSRAM memory, or they offload the work to a server. Nobody had characterized the bare classic ESP32: no PSRAM, only about a hundred and fifty kilobytes of usable memory. It is a three-dollar board, and by Espressif's own ranking it is the weakest WiFi sensor in the family. Yet that is exactly the chip a real, cost-sensitive product would choose. So my question was: how far can it actually go?

\slide{6}{Approach: the deployability frontier}{$\approx$40s}
\cue{Three model tiers and the train to flash pipeline.}
My approach was to build what I call a deployability frontier. I took three tiers of models: classical learners like decision trees, tiny neural networks, and five standard deep architectures. I ran all of them across two public datasets, each with seven activities. For every single model I asked one honest question: does it actually convert, fit, and run on the bare chip? And importantly, I tested the CSI-HAR dataset leave-one-user-out, which means the model is always tested on a person it never saw during training.

\slide{7}{Finding 1: deployment is not the wall}{$\approx$45s}
\cue{Arena sizes (10 to 54 kB) and the memory chart.}
My first finding genuinely surprised me. I expected the deep models to be too big, but they are not. I flashed them to the real board, and a full deep CNN, and even a Transformer, both run on the bare ESP32, using only ten to fifty-four kilobytes. The trick is that the weights sit in the four-megabyte flash, and only the activations need the scarce SRAM. So there is no general memory wall. The chip is far more capable than people assume.

\slide{8}{The real barriers}{$\approx$35s}
\cue{Table: CNN and Transformer fit; GRU, SSM, KAN do not.}
So what does fail, and why? Two narrow things, and neither is about memory size. First, convertibility: the recurrent and state-space models compile to operators that TensorFlow Lite Micro simply does not support, so they never reach the chip. Second, the Kolmogorov-Arnold network converts but fails to allocate at runtime. Everything else, the CNN, the Transformer, the tiny models, and the classical ones, fits and runs.

\slide{9}{The generalization gap (headline)}{$\approx$50s}
\cue{The signal dropping from 96\% to 63\%, the ``-29 points'' marker.}
And this is the most important slide of the talk. If deployment is not the wall, then what is? It is generalization. Watch what happens to the same model when I change only one thing: how the data is split. On a random split, which lets the model memorize each person, it scores ninety-six percent. But test it on a person it has never seen, and it falls to sixty-three percent. That is a twenty-nine-point collapse, and it happens to every model, regardless of size. A random split basically lets the model cheat by memorizing each user's body signature. So the real open problem is not fitting the model on the chip; it is recognizing someone new.

\slide{10}{Results at a glance}{$\approx$30s}
\cue{Stat cards and the accuracy chart.}
To put numbers on it: a tiny CNN reaches about ninety-seven percent on the saturated benchmark, using only eight to thirteen kilobytes of RAM, and running in fourteen to a hundred and seventeen milliseconds on the real board. The classical Random Forest runs in about forty-five microseconds. So the small models are not a compromise; they match the deep models at a tiny fraction of the cost.

\slide{11}{Takeaways}{$\approx$35s}
\cue{Four numbered takeaways.}
So, four takeaways. One: deployment is solved; even deep models fit the cheapest ESP32. Two: generalization is the real wall, with about a twenty-nine-point drop on a new person. Three: we should report subject-independent accuracy, because random-split numbers are misleading. And four: small models are the smart default for this hardware.

\slide{12}{Thank you / Q\&A}{$\approx$15s}
\cue{Closing slide with your details.}
To sum up: the cheapest WiFi chip can run this AI just fine. The next frontier is helping it recognize a person it has never seen before. Thank you very much for listening. I would be happy to take any questions.

\vspace{10pt}
\hrule
\vspace{6pt}
{\small\bfseries\color{teal} Likely questions and short answers}
\begin{itemize}[leftmargin=1.4em,itemsep=2pt]
  \item \textbf{Why is CSI-HAR accuracy low (around 63\%)?} Because it is tested on a completely unseen person. That is the honest, real-world number; random splits hide it.
  \item \textbf{Did you test on real WiFi live?} I used two public datasets measured on real hardware. The on-device latency and memory are measured on my physical ESP32.
  \item \textbf{How would you close the generalization gap?} A short per-user calibration, or cross-subject methods like domain adaptation, rather than a bigger model.
  \item \textbf{Why does the KAN fail?} It converts to int8 but the default runtime cannot allocate its quantized polynomial layers; it is a tooling limit, not a fundamental one.
\end{itemize}

\end{document}
